\scanpage{384}
\mspage{300}
% archive sheet number: 383

\deleted{\unsure{on aura}}\unsure{pour} tout $r$, comme
\[
I_r\overset{\mathrm{def}}{=}\unsure{\varphi|\phi}^{-1}(r)\subset I_{\alpha_\pi}
\]
est cofinal dans $I_{\alpha_\pi}$, on aura donc
\[
j=\varinjlim_{\alpha\in\deleted{\unsure{\varphi^{-1}(r)}}I_r}i_\alpha,
\]
et \deleted{comme les} en vertu de $(*)$ les
$i_\alpha$ pour $\alpha\in I_r$ sont dans $J_r$, donc
par passage à la limite, utilisant
l'hypothèse c), on trouve que $j\in J_r$,
donc \deleted{$j\in$}
\[
j\in\bigcap_{r\in A}J_r=J,
\]
et comme on choisira
\[
i_0=i,\tag*{$(**)$}
\]
on aura bien trouvé $i\longrightarrow j$ avec $i\in J$.
Donc il suffit de construire le système
inductif $(i_\alpha)$ satisfaisant $(*)$ et $(**)$. 

Si \inserted{le syst\textsuperscript{ème} inductif est construit jusque \unsure{compris|et y compris}} $i_\alpha\in\operatorname{Ob}J_{\varphi(\alpha)}$, \deleted{\unsure{est cons}} on continue
par
\[
i_{\alpha\deleted{\unsure{+1}}}\longrightarrow i_{\alpha+1}\in\operatorname{Ob}J_{\varphi(\alpha+1)},
\]
ce qui est possible puisque $J_{\varphi(\alpha+1)}$ est
